\documentclass{article}

\usepackage[education, final]{neurips_2026}

\usepackage[utf8]{inputenc}
\usepackage[T1]{fontenc}
\usepackage{hyperref}
\usepackage{url}
\usepackage{booktabs}
\usepackage{amsmath}
\usepackage{xcolor}
\usepackage{graphicx}
\setlength{\textfloatsep}{8pt plus 2pt minus 2pt}
\setlength{\abovecaptionskip}{4pt}

\title{Speculative Decoding: What's Measured, What's Missed, and What's Next}

\author{%
  Lily Zhang \and Madison Kanna
}

\begin{document}

\maketitle

\begin{abstract}
An interactive, plain-language resource that teaches speculative decoding as the field evaluates it, then examines the evaluation itself. Learners trace three generations of draft models (EAGLE-3, DFlash, DSpark), learn to read \emph{acceptance length} correctly, and see why rejection-sampling verification makes the base method provably lossless. The standard harness never grades outputs, the theorem covers only exact verification, and the measured behavior gap between owner-trained and vendor-assembled acceleration reaches 5.6 points on long-tail domains. A one-GPU, one-afternoon lab lets learners reproduce acceptance numbers, sweep the confidence threshold that converts lossless into lossy, and grade the outputs themselves.
\end{abstract}

\section{Concept Summary}
The market wants more tokens. Agent workflows chain dozens of model calls per task, reasoning models spend thousands of tokens thinking before they answer, and every product built on either one pays for latency twice, in compute and in user patience. The ideal answer is lossless inference acceleration: serve the same model faster without changing what it outputs. That demand has pulled acceleration steadily deeper into the model itself. Draft models for speculative decoding evolved through three generations in two years: \textbf{EAGLE-3} reuses target hidden states for feature-level drafting ($\sim$3.9 tokens/pass); \textbf{DFlash} drafts a whole block in one diffusion pass ($\sim$4.4); \textbf{DSpark} adds a sequential head and confidence-scheduled trimming ($\sim$5.1; Qwen3-4B code average, DSpark Table~1). Frontier labs moved in the same direction from the training side: DeepSeek pushed FP8 into pre-training, gpt-oss shipped MXFP4 with quantization-aware training, and Kimi K3 ships the speculator itself, fine-tuned and validated before release. The \href{https://lily-ideas.vercel.app/neurips-education/}{interactive walkthrough} carries one worked example across all four decoding stages; in the DSpark stage a low-confidence draft token is trimmed before verification and corrected by the target, which is where lossless ends.

Inference providers now run some combination of quantization, speculative decoding, and serving-level optimizations on nearly every hosted model, and the field reports progress in tokens per second while describing the acceleration as lossless. What no one owns is the quality of the accelerated output. \emph{What's measured:} a small draft proposes $K$ tokens and the target verifies all $K$ in a single forward pass at nearly the cost of one, and rejection-sampling verification guarantees the output distribution exactly, so evaluation collapsed onto acceptance length and tokens/s; draft training converged onto the same number (the LK loss is the negative log of per-token acceptance probability). \emph{What's missed:} in DeepSpec's evaluation harness the nine standard datasets supply only prompts, with no grader or accuracy metric anywhere in the code, so the lossless claim rests entirely on a theorem that covers exact verification only; quantization, KV routing, disaggregation, and relaxed acceptance sit outside it, and the measured behavior gap between owner-trained and vendor-assembled acceleration reaches 5.6 points on long-tail domains, while verification concentrates in low-entropy math and code, exactly where acceleration is least likely to misbehave. \emph{What's next:} no published objective trains a draft to preserve the target's behavior on the long tail.

\section{Prerequisites and Technical Level}
\textbf{Level:} introductory-to-intermediate; accessible to any engineer or student who knows that an LLM generates text one token at a time. \textbf{Prerequisites:} basic familiarity with transformer inference (a forward pass produces a next-token distribution); no RL, training, or GPU-kernel background required. The rejection-sampling losslessness argument is presented intuitively, with the accept/repair rule optional for readers who want the math.

\section{Learning Objectives and Outcomes}
After engaging with this resource, a learner can:
\begin{itemize}
  \item explain \emph{why} speculative decoding is possible (verifying $K$ tokens $\approx$ the cost of one), and read \emph{acceptance length} correctly: accepted proposals per verification step, not saved passes;
  \item trace the progression from EAGLE-3 to DFlash to DSpark, name the bottleneck each stage removes, and explain why training objective and evaluation metric converged onto the same number;
  \item \textbf{demonstrate hands-on} that acceptance rate is not accuracy: sweep the confidence threshold in the official evaluation code, grade the outputs, and watch the two numbers move in opposite directions;
  \item articulate the domain mismatch, verification concentrates where acceptance is naturally highest (math, code) while usage concentrates elsewhere, and pose the open question: what would a draft-training objective that preserves long-tail behavior look like?
\end{itemize}

\section{Grounding in NeurIPS Research (2022--2026)}
The three generations the resource teaches are themselves recent papers at NeurIPS and its sister venues:
\begin{itemize}\small
  \item \href{https://openreview.net/forum?id=4exx1hUffq}{\textbf{EAGLE-3}} (NeurIPS 2025): feature-level drafting with multi-layer fusion and a training-time test, the current autoregressive-draft baseline.
  \item \href{https://icml.cc/virtual/2026/poster/64301}{\textbf{DFlash}} (ICML 2026): a block-diffusion drafter that generates the whole block in one forward pass, over $6\times$ lossless acceleration and up to $2.5\times$ over EAGLE-3.
  \item \href{https://arxiv.org/abs/2607.05147}{\textbf{DSpark}} (DeepSeek, 2026): confidence-scheduled verification on the block-parallel backbone, serving DeepSeek-V4-Flash 60--85\% faster in production.
\end{itemize}
The walkthrough places them in the NeurIPS speculative-decoding line they extend, SpecTr (NeurIPS 2023), Sequoia (NeurIPS 2024), and SpecExec (NeurIPS 2024), alongside the foundational papers (Leviathan et al., 2023; Chen et al., 2023), FP8 pre-training in DeepSeek-V3 (2024), MXFP4 QAT in gpt-oss (2025), Judge Decoding (ICLR 2025), and Kimi K3 shipping the speculator as part of post-training (2026).

\section{Teaching Materials Summary}
All materials are original and created for this track: (1) an \textbf{interactive self-contained HTML article} following the What's Measured / What's Missed / What's Next arc, with a walkthrough with adjustable draft length; (2) a \textbf{one-afternoon lab} (single GPU): reproduce the reported EAGLE-3 acceptance length (2.4--2.8) on Qwen3-8B, compare the three released DeepSpec checkpoints on one prompt set, sweep \texttt{--confidence-threshold} while grading outputs for correctness, and a pick-a-domain exercise outside math and code; (3) a \textbf{video walkthrough} recorded from the site

\end{document}
